\fsection{Especificaciones generales}
\fsubsection{Material}
Para el desarrollo del proyecto se utilizará el siguiente equipamiento:
\begin{itemize}
	\item Cámara de profundidad 3D Orbbec Femto Mega (imagen en color y nube de puntos).
	\item PC industrial Siemens SIMATIC IPC BX-59A con el software SIMATIC Robot Pick AI.
	\item PLC Siemens S7-1500, CPU 1516-3 PN/DP (firmware V2.9).
	\item Robot colaborativo Universal Robots UR3e (6 ejes).
	\item Pinza de vacío eléctrica OnRobot VGC10 como efector final (vacío generado por bomba integrada, sin aire comprimido).
	\item PC de ingeniería con una IP en el rango 192.168.15.xxx, desde el que se programa el PLC y se accede
	      al frontend web del IPC.
	\item Cableado de red Profinet.
	\item Patrón de calibración impreso en folio, para calibrar la cámara y el robot.
	\item Mesa de trabajo móvil (con ruedas) de 1400 × 800 mm, sobre la que se montan el robot, el contenedor y
	      las bandejas.
	\item Contenedor de piezas de 345 × 270 × 60 mm (interior de 325 × 250 × 55 mm).
	\item Bandeja de colocación para las piezas recogidas, más otras dos reservadas para la ampliación de
	      clasificación.
\end{itemize}

\fsubsection{Software}
\begin{itemize}
	\item TIA Portal V18 o superior de Siemens, para la programación del PLC a partir del proyecto de la nota de aplicación
	      \textit{Robot Pick AI Pro Machine Vision Software for AI-based Universal Bin Picking}.
	\item SIMATIC Robot Pick AI (contenedores Docker en el IPC: drivers de la cámara, frontend web de
	      configuración y modelo de IA preentrenado).
	\item Orbbec Viewer, para dejar la cámara con la configuración que espera el software del IPC.
	\item PolyScope de Universal Robots, para el programa (\texttt{.urp}) y el archivo de instalación del robot.
\end{itemize}
